\documentclass[11pt]{article}
\usepackage[margin=1in]{geometry}
\usepackage{amsmath,amssymb}
\usepackage[hidelinks]{hyperref}

\begin{document}

\section*{Human Review: Asymptotic Dimension Zero ACL Counterexample}

\noindent Original automatic proof: \texttt{solution\_packet.pdf}

\noindent Checked date: 19 June 2026

\noindent Status: Manually checked.

\noindent Verdict: The proof appears correct.

\section*{Human Comments}

The proposed construction gives a correct negative answer to Problem 2.15 of
Braga--Lancien, \emph{Asymptotic coarse Lipschitz equivalence}, as the problem is
stated for arbitrary metric spaces.  The source question asks whether two
asymptotically coarse Lipschitz equivalent metric spaces of asymptotic dimension
zero must be coarse, or coarse Lipschitz, equivalent.  The packet constructs
spaces \(X\) and \(Y\) with
\[
  \operatorname{asdim}(X)=\operatorname{asdim}(Y)=0,
\]
which are asymptotically coarse Lipschitz equivalent but are not even coarsely
equivalent.  Since coarse Lipschitz equivalence implies coarse equivalence, this
is the stronger negative conclusion.

The idea is clear.  The space \(X=\{x_n:n\in\mathbb N\}\) is a sparse sequence
with
\[
  d_X(x_n,x_m)=2^n+2^m\qquad(n\neq m).
\]
The space \(Y\) replaces each point \(x_n\) by a fiber \(\{n\}\times A\), with
internal distance \(n\), while keeping different fibers at distance
\(2^n+2^m\).  Collapsing each fiber gives a coarse Lipschitz map
\(g:Y\to X\), and choosing one representative in each fiber gives a coarse
Lipschitz map \(f:X\to Y\).  One has \(gf=\operatorname{Id}_X\), while
\(fg\) moves a point in the \(n\)-th fiber by at most \(n\).  Since this is
sublinear compared with the distance scale \(2^n\), the maps witness asymptotic
coarse Lipschitz equivalence.

The asymptotic dimension zero verification is also sound.  At a fixed scale
\(r\), the finitely many small fibers may be grouped into one bounded set, while
for all sufficiently large \(n\), the points inside the \(n\)-th fiber are already
farther than \(r\) apart, so singletons handle the large fibers.

The proof that the spaces are not coarsely equivalent is correct.  If a coarse
equivalence existed, there would be coarse maps \(F:X\to Y\) and \(G:Y\to X\)
with \(FG\) uniformly close to \(\operatorname{Id}_Y\).  Thus there is
\(C<\infty\) such that
\[
  d_Y(FG(y),y)\le C\qquad(y\in Y).
\]
Choosing a fiber whose internal distance is larger than \(C\) forces
\(FG\) to be the identity on that fiber.  In the packet's original version,
where the fiber is uncountable and \(X\) is countable, this gives an immediate
cardinality contradiction.

\section*{Suggested Countable Variant}

The uncountability of the set \(A\) is not actually needed.  A more natural
version of the example is obtained by taking
\[
  A=\mathbb N,
\]
so that both \(X\) and \(Y=\mathbb N\times\mathbb N\) are countable.  The
asymptotic coarse Lipschitz equivalence and the asymptotic-dimension-zero
argument are unchanged.

Only the final non-coarse-equivalence argument needs a small adjustment.  If
\(F:X\to Y\) and \(G:Y\to X\) were coarse maps with
\(FG\sim\operatorname{Id}_Y\), let \(C\) be the closeness constant and choose
\(n>C\).  Then \(FG\) must be the identity on the infinite fiber
\[
  Y_n=\{n\}\times\mathbb N.
\]
Hence
\[
  F(G(Y_n))=Y_n.
\]
But \(Y_n\) has finite diameter \(n\).  Since \(G\) is coarse, \(G(Y_n)\) has
bounded diameter in \(X\).  Every bounded subset of the sparse space \(X\) is
finite, so \(G(Y_n)\) is finite, and therefore \(F(G(Y_n))\) is finite.  This
contradicts the infinitude of \(Y_n\).

Thus the example can be made countable.  It still does not have bounded
geometry or locally finite fibers, so it should still be viewed as answering the
literal formulation of the source problem rather than any stronger bounded-
geometry version.

\section*{Review Outcome}

Archive this packet as a correct literal-wording counterexample to Problem 2.15.
For presentation, I recommend replacing the uncountable fiber set by the
countable choice \(A=\mathbb N\), together with the adjusted finite-image
argument above.  This makes the example more natural while preserving the same
obstruction.

\end{document}
